%% kapa: exam=1 round=23 year=2012 period=2 form=A subject=appraisal_law subject_ko=부동산관계법규 seq=3 number=43
%% source: Q-Net 감정평가사 기출문제 (cst00309) artlSeq=5052580 file=raw/1008/5052580/2012년 제23회 감정평가사 제1차시험 기출문제(2교시 A형).pdf
%% conversion: kordoc 4.24.51 -> tools/gen1.py (auto); review=proofread-sonnet; flags=none
\begin{question}{43}
국토의 계획 및 이용에 관한 법령상 용도지역에 관한 설명이다. ( )에 들어갈 용어가 옳게 연결된 것은?
\begin{framedtext}
( ㉠ ) : 공동주택 중심의 양호한 주거환경을 보호하기 위하여 필요한 지역 \par
( ㉡ ) : 근린지역에서의 일용품 및 서비스의 공급을 위하여 필요한 지역 \par
( ㉢ ) : 주로 농업적 생산을 위하여 개발을 유보할 필요가 있는 지역 \par
\end{framedtext}
\noindent\hspace*{2em}\makebox[9.5em][l]{㉠}\makebox[7em][l]{㉡}㉢ \par
\begin{choices}
  \item \makebox[9.5em][l]{제2종일반주거지역}\makebox[7em][l]{근린상업지역}생산녹지지역
  \item \makebox[9.5em][l]{제2종일반주거지역}\makebox[7em][l]{근린상업지역}농림지역
  \item \makebox[9.5em][l]{제2종일반주거지역}\makebox[7em][l]{일반상업지역}농림지역
  \item \makebox[9.5em][l]{제2종전용주거지역}\makebox[7em][l]{근린상업지역}생산녹지지역
  \item \makebox[9.5em][l]{제2종전용주거지역}\makebox[7em][l]{일반상업지역}생산녹지지역
\end{choices}
\end{question}
\answer{4}
